\documentclass[11pt,a4paper]{article}
\usepackage[margin=1in]{geometry}
\usepackage{amsmath,amssymb}
\usepackage{graphicx}
\usepackage{float}
\usepackage{enumitem}
\usepackage{microtype}
\usepackage{caption}
\usepackage{tabularx}
\usepackage{hyperref}

\captionsetup{font=small,labelfont=bf,skip=6pt}

\setlist[itemize]{leftmargin=1.2cm}
\setlist[enumerate]{leftmargin=1.2cm}

\newcommand{\problembox}[3]{%
  \par\addvspace{0.42em}%
  \noindent
  \setlength{\fboxsep}{7.5pt}%
  \setlength{\fboxrule}{0.95pt}%
  \fbox{%
    \begin{minipage}{\dimexpr\linewidth-2\fboxsep-2\fboxrule\relax}
    \begin{tabularx}{\linewidth}{@{}>{\raggedright\arraybackslash\bfseries}p{2.1em} @{\hspace{0.9em}} >{\raggedright\arraybackslash}X @{\hspace{0.9em}} >{\raggedleft\arraybackslash\bfseries}p{2.4em}@{}}
    #1 & #2 & #3
    \end{tabularx}
    \end{minipage}%
  }%
  \par\addvspace{0.32em}%
}

\title{\textbf{Thermodynamics of ferroelectrics}}
\author{\normalsize Y24 (2024) --- English translation}
\date{}
\begin{document}
\maketitle

Ferroelectrics are substances that, at a sufficiently low temperature, possess a non-zero electric polarization in the absence of an external electric field. In this respect, they are analogous to ferromagnets. This problem considers the thermodynamic properties of ferroelectrics, in particular the electrocaloric effect (temperature change due to electric field variation), which allows for the use of ferroelectric films for surface cooling.

\begin{figure}[H]
    \centering
    \includegraphics[width=0.72\textwidth]{T4_EN_assets/fig1.jpg}
    \caption{Surface cooling using the electrocaloric effect in a ferroelectric film.}
\end{figure}

The state of a ferroelectric is characterized by polarization $P$ (dipole moment per unit volume), electric induction $D$, and electric field strength $E$. For a flat sample with uniform polarization, we treat these as scalar projections onto the axis of polarization. The volumetric density of electrical energy changes as $dw = E dD$. For a constant volume $V$, the internal energy change is $dU = \delta Q + V E dD$, where $\delta Q$ is the heat supplied.

\textbf{Mathematical Supplement:} Consider a function $f(x,y)$. The partial derivative with respect to $x$ at constant $y$ is denoted as $\left( \frac{\partial f}{\partial x}\right)_y$. For a triple of variables $x, y, z$ related by an equation of state, the following identity holds:
$$ \left( \frac{\partial x}{\partial y}\right)_z \left( \frac{\partial y}{\partial z}\right)_x \left( \frac{\partial z}{\partial x}\right)_y = -1. \quad(1) $$
Additionally, the reciprocal rule states $\left( \frac{\partial x}{\partial y}\right)_z = 1 / \left( \frac{\partial y}{\partial x}\right)_z$. For entropy $S$ and temperature $T$, the relationship $(\partial S / \partial V)_T \neq (\partial S / \partial V)_p$ generally holds.

\section*{Part A. Electrocaloric effect in ferroelectrics (5.5 points)}

\problembox{A1}{Consider one mole of an ideal gas. Using its pressure $p$, volume $V$, and temperature $T$, verify identity (1).}{0.50}

\problembox{A2}{Determine the total heat received by a ferroelectric per cycle performed in coordinates $(E, P)$. Express the answer in terms of $E_1, E_2, P_1, P_2, V$.}{0.70}

\begin{figure}[H]
    \centering
    \includegraphics[width=0.45\textwidth]{T4_EN_assets/fig2.jpg}
    \caption{Rectangular cycle in $(E, P)$ coordinates.}
\end{figure}

\problembox{A3}{Consider the infinitesimal parallelogram cycle $ABCD$. Indicate the direction of traversal for which the heat received is positive. Justify your answer.}{0.40}

\begin{figure}[H]
    \centering
    \includegraphics[width=0.45\textwidth]{T4_EN_assets/fig3.jpg}
    \caption{Infinitesimal parallelogram cycle for thermodynamic analysis.}
\end{figure}

The following points prove the identity $\left( \frac{\partial S}{\partial E} \right)_T = \left( \frac{\partial P}{\partial T}\right)_E \quad(2)$, where $S$ is the entropy per unit volume.

\problembox{A4}{Let $BC$ and $AD$ be isotherms at $T_1$ and $T_2$ ($|T_1 - T_2| \ll T_1$). Express the total heat $Q$ for the cycle in terms of $V, (\partial P/\partial T)_E, T_1, T_2$ and $dE = E_D - E_A \approx E_C - E_B$.}{0.50}

\problembox{A5}{Find the heat $Q_+$ supplied to the ferroelectric. Express it in terms of $V, T_1, T_2, \Delta E$ and $(\partial S/\partial E)_T$.}{0.50}

\problembox{A6}{Using the efficiency of a Carnot cycle and the results of A4--A5, prove relation (2). Assume the total heat received equals the useful work.}{0.30}

\problembox{A7}{Now let $AB$ and $CD$ be isotherms. Express $Q$ in terms of $V, dE = E_B - E_A, T_1, T_2, (\partial E/\partial T)_P$ and $(\partial P/\partial E)_T$.}{0.60}

\problembox{A8}{Using identity (1) for $E, P, T$, express $Q$ from A7 in terms of $V, T_1, T_2, dE$ and $(\partial P/\partial T)_E$.}{0.50}

\problembox{A9}{Write the expression for $Q_+$ supplied in this cycle using $V, T_1, T_2, \Delta E$ and $(\partial S/\partial E)_T$.}{0.50}

\problembox{A10}{Using the Carnot efficiency and the definition of useful work, prove formula (2) for this case.}{0.30}

\problembox{A11}{Using (2) and relation (1) for $S, T, E$, obtain: $\left( \frac{\partial T}{\partial E}\right)_S = -\frac{T}{C_E} \left(\frac{\partial P}{\partial T} \right)_E$, where $C_E = \left( \frac{\delta Q}{dT}\right)_E$ is the specific heat capacity at constant field.}{0.70}

\section*{Part B. Microscopic model of a ferroelectric (4.5 points)}

Consider ions in an effective potential:
$$ U(x) = U_0 \left( \frac{a_2}{2} \left( \frac{x}{x_0}\right)^2 - \frac{a_4}{4} \left( \frac{x}{x_0}\right)^4 + \frac{a_6}{6} \left( \frac{x}{x_0}\right)^6\right) $$
where $U_0, x_0, a_4, a_6 > 0$ are constants, and $a_2 = \alpha (T - T_c)/T_c$. Let $n$ be the ion concentration and $q$ the ion charge.

\problembox{B1}{Find the equilibrium positions. Express them through $x_0, a_2, a_4, a_6$. Determine the number of equilibria based on $a_2$ and indicate their stability.}{1.40}

\problembox{B2}{For $a_4=1, a_6=1$, numerically determine the value $a_{2c}$ at which a non-zero $x$ becomes the global energy minimum.}{0.70}

\problembox{B3}{Find the spontaneous polarization $P(T)$. Express it using $n, q, x_0, a_4, a_6, \alpha$ and $t = (T-T_c)/T_c$. Assume all ions shift in the same direction.}{0.40}

\problembox{B4}{Calculate the derivative $(\partial P/\partial T)_{E=0}$. Express it through $n, q, x_0, a_4, a_6, \alpha, t, T_c$. Plot this derivative vs. $t$.}{0.70}

\problembox{B5}{Apply a weak field $E$. Write the equilibrium equation and find the polarizability $\chi = \left(\frac{\partial P}{\partial E} \right)_{T, E=0}$. Express it using $n, q, x_0, U_0, a_4, a_6, \alpha, t$ and the equilibrium position $x$.}{1.30}

\vspace{1em}
\noindent\textit{Source:} https://pho.rs/p/3851

\end{document}